\par\noindent\textbf{(1)} 两质点的位置线性近似为 \(x_1=l\theta_1,x_2=l(\theta_1+\theta_2)\)，故 \(T=\tfrac12 ml^2(2\dot\theta_1^2+2\dot\theta_1\dot\theta_2+\dot\theta_2^2)\)，\(V=\tfrac12 mgl(2\theta_1^2+\theta_2^2)\)。运动学交叉项不能忽略。
\par\noindent\textbf{(2)} 质量矩阵 \(M=ml^2\begin{pmatrix}2&1\\1&1\end{pmatrix}\)，刚度矩阵 \(K=mgl\operatorname{diag}(2,1)\)。令 \(\theta=ve^{i\omega t}\)，\(\det(K-\omega^2M)=0\) 给 \(\omega_{\pm}^2=(g/l)(2\pm\sqrt2)\)，两根均正。
\par\noindent\textbf{(3)} 从首行 \([2g-2l\omega^2]v_1-l\omega^2v_2=0\)，得 \(v_2/v_1=2g/(l\omega^2)-2\)。低频模得到 \(v_-=(1,\sqrt2)^T\)，高频模为 \(v_+=(1,-\sqrt2)^T\)，分别表现同相与反相运动。
\par\noindent\textbf{(4)} 两模态必须满足 \(v_-^TMv_+=0\) 而不是普通欧氏内积正交。代入得 \(ml^2[2+(-\sqrt2+\sqrt2)-2]=0\)。因此任意小扰动可分解为两种独立的线性正常振动，其能量在模态坐标中解耦。
